% Source: https://artofproblemsolving.com/community/c3h1476946_plane_traveling_with_wind_quadratic_mao_1990
% Stable identity: aops-topic-1476946-q01
With no wind the plane covers $2000$ km in $10$ hours, so its airspeed is $200$ km/h. With the wind it flies from $A$ to $B$ at $200+k$ and back at $200-k$, and the total time is $10\dfrac{5}{12}=\dfrac{125}{12}$ hours:
\[
\dfrac{1000}{200+k}+\dfrac{1000}{200-k}=\dfrac{1000\cdot400}{40000-k^2}=\dfrac{125}{12}.
\]
Instead of expanding, cross-multiply and divide by $125$:
\[
40000-k^2=\dfrac{400000\cdot12}{125}=38400,\qquad k^2=1600,\qquad k=\boxed{40}.
\]
The key simplification is to combine the two fractions first: their numerators add up to $1000\cdot400$ and the denominator is a difference of squares, so the ``complicated quadratic'' becomes $k^2=1600$.
